\documentclass[11pt]{article}
\usepackage[T1]{fontenc}
\usepackage{textcomp}
\usepackage[margin=0.75in]{geometry}
\usepackage{amsmath,amssymb,amsthm}
\usepackage{array,booktabs}
\usepackage{hyperref}
\setlength{\parindent}{0pt}
\newtheorem{theorem}{Theorem}
\theoremstyle{definition}
\newtheorem{definition}{Definition}
\title{Flow Proof Artifact: prop-06-same-ratio-to-third}
\author{Generated by \texttt{flow doc proof}}
\date{Flow Proof Book}
\begin{document}
\maketitle
Magnitudes which have the same ratio to the same magnitude are equal in ratio to one another.\\[0.5em]
\textbf{Source.} euclid — Elements, Book V, Proposition 6\\[0.5em]
\bigskip
\noindent{\large \textbf{Derived fact 1.} \textit{Magnitudes which have the same ratio to the same magnitude are equal in ratio to one another} \hfill the Euclidean plane \cdot Euclid Book V \cdot Proposition 6: magnitudes with the same ratio to a third are proportional to one another}\par
\smallskip\noindent\textit{Coordinate: the Euclidean plane \cdot Euclid Book V \cdot Proposition 6: magnitudes with the same ratio to a third are proportional to one another}\par
\smallskip\noindent\textit{Source: euclid — Elements, Book V, Proposition 6}\par
\smallskip\noindent\textit{Built on: proposition 4: equimultiples preserve proportion, for Euclid Book V on the Euclidean plane}\par
\medskip
\noindent\textbf{Goal.} Magnitudes which have the same ratio to the same magnitude are equal in ratio to one another\par
\noindent$\displaystyle \text{same ratio to third implies equal ratios}$\par
\medskip
\noindent
\begin{tabular}{@{} >{\raggedright\arraybackslash}p{0.06\textwidth} >{\raggedright\arraybackslash}p{0.47\textwidth} >{\raggedleft\arraybackslash}p{0.41\textwidth} @{}}
\textbf{\#} & \textbf{Proof} & \textbf{Mathematics} \\
\midrule
\textbf{1.} & We prove that proposition 6: magnitudes with the same ratio to a third are proportional to one another for Euclid Book V the Euclidean plane. & \\[0.45em]
\textbf{2.} & Let a = first. & \\[0.45em]
\textbf{3.} & Let b = second. & \\[0.45em]
\textbf{4.} & Let c = common\_third. & \\[0.45em]
\textbf{5.} & From step 2, step 3, and step 4, we can deduce that ratio(a, b) equals ratio(d, e). & $\displaystyle ratio(a, b) = ratio(d, e)$ \\[0.45em]
\textbf{6.} & We invoke the derived fact governing Euclid Book V the Euclidean plane: proposition 4: equimultiples preserve proportion, for Euclid Book V on the Euclidean plane. & \\[0.45em]
\textbf{7.} & From step 6, this implies same ratio to third implies equal ratios. Hence proven. & $\displaystyle \text{same ratio to third implies equal ratios}$ \\[0.45em]
\bottomrule
\end{tabular}
\label{thm:Geometry-Euclid Book V-proposition_6_magnitudes_with_the_same_ratio_to_a_third_are_proportional_to_one_another}
\par\medskip\hrule
\end{document}
